\paragraph{Các quan hệ chính}
\mbox{}\par

Một tài khoản có thể gửi hoặc nhận nhiều yêu cầu kết bạn, xuất hiện trong
nhiều quan hệ bạn bè và quan hệ chặn, đồng thời tham gia nhiều hội thoại.
Với các quan hệ có hai vai trò, hai thuộc tính cùng tham chiếu đến
\texttt{UserAccount.id} nhưng có ý nghĩa khác nhau, chẳng hạn người gửi
và người nhận yêu cầu kết bạn.

\begin{itemize}
  \item \textbf{Tài khoản và quan hệ bạn bè:}
    \texttt{FriendRequest.senderId} và \texttt{receiverId},
    hai thuộc tính của khóa chính \texttt{Friendship},
    cùng \texttt{UserBlock.blockerId} và \texttt{blockedUserId}
    đều tham chiếu đến tài khoản tương ứng.
  \item \textbf{Tài khoản và hội thoại:}
    \texttt{Conversation.creatorId} tham chiếu người tạo.
    Quan hệ nhiều--nhiều giữa tài khoản và hội thoại được tách qua
    \texttt{ConversationMember}, có khóa chính kép
    \texttt{(conversationId, userId)}.
  \item \textbf{Hội thoại và tin nhắn:}
    mỗi tin nhắn thuộc một hội thoại. Cặp
    \texttt{(conversationId, senderId)} trong \texttt{Message}
    tham chiếu khóa chính của \texttt{ConversationMember}, giúp bảo đảm
    người gửi có bản ghi thành viên trong đúng hội thoại.
    Điều kiện còn đang tham gia tại thời điểm gửi cần được backend kiểm tra thêm.
  \item \textbf{Tin nhắn và lịch sử cá nhân:}
    \texttt{HiddenMessage} liên kết một người dùng với một tin nhắn.
    Khóa chính kép \texttt{(userId, messageId)} ngăn ghi nhận trùng
    cùng một lần ẩn tin của người dùng đó.
  \item \textbf{Tin nhắn và báo cáo spam:}
    \texttt{SpamReport} tham chiếu người báo cáo, tài khoản bị báo cáo
    và tin nhắn liên quan; một tin nhắn có thể xuất hiện trong nhiều báo cáo.
  \item \textbf{Tài khoản và hoạt động:}
    mỗi bản ghi \texttt{LoginHistory} hoặc \texttt{AppOpenEvent}
    thuộc một tài khoản; một tài khoản có thể có nhiều bản ghi theo thời gian.
\end{itemize}

\paragraph{Ràng buộc và quy tắc dữ liệu dự kiến}
\mbox{}\par

Các khóa chính, khóa ngoại và ràng buộc duy nhất được dự kiến thể hiện trong
CSDL. Những quy tắc liên quan đến quyền truy cập, trạng thái tham gia hoặc
thay đổi đồng thời nhiều bản ghi cần được kiểm tra thêm tại backend;
không xem các khóa trong ERD là đủ để thực hiện toàn bộ nghiệp vụ.

\begin{enumerate}
  \item \textbf{Tài khoản:} tên đăng nhập và email là duy nhất;
    mật khẩu được lưu dưới dạng đã băm. \texttt{systemRole} phân biệt
    Admin hệ thống với User; \texttt{status} phân biệt tài khoản đang hoạt động
    với tài khoản bị khóa.
  \item \textbf{Bạn bè và chặn:} hai người trong một quan hệ phải khác nhau.
    \texttt{Friendship} lưu cặp ID theo quy ước
    \texttt{userId1 < userId2}, kết hợp khóa chính kép để không lưu
    hai chiều của cùng một quan hệ. Quan hệ chặn có hướng nên A chặn B
    không đồng nghĩa B chặn A.
  \item \textbf{Hội thoại và thành viên:} \texttt{type} nhận
    \texttt{PRIVATE} hoặc \texttt{GROUP}; tên nhóm được dùng cho hội thoại nhóm.
    Hội thoại riêng có hai người tham gia; khi tạo nhóm cần người tạo và
    ít nhất một người dùng khác. Khóa chính kép của thành viên ngăn một
    người xuất hiện hai lần trong cùng hội thoại.
  \item \textbf{Quyền trong nhóm:} \texttt{memberRole} nhận
    \texttt{ADMIN} hoặc \texttt{MEMBER}, độc lập với quyền quản trị hệ thống.
    Chỉ admin nhóm được xóa thành viên theo yêu cầu đề bài.
    Khi gửi tin, backend cần kiểm tra trạng thái thành viên và quyền hiện tại.
  \item \textbf{Báo cáo spam:} người báo cáo phải có quyền xem nội dung
    liên quan; tin nhắn được báo cáo phải thuộc tài khoản bị báo cáo.
    Backend cần kiểm tra các điều kiện này ngoài việc kiểm tra các khóa ngoại.
  \item \textbf{Tính đồng bộ:} khi chấp nhận kết bạn, cập nhật yêu cầu và tạo
    quan hệ bạn bè trong cùng giao dịch. Khi chặn và hủy kết bạn, các thay đổi
    liên quan cũng cần được xử lý đồng bộ để tránh trạng thái mâu thuẫn.
\end{enumerate}

\paragraph{Xóa lịch sử và mã hóa nhóm}
\mbox{}\par

Thiết kế đề xuất hiểu thao tác xóa lịch sử là xóa khỏi phía người thực hiện.
\texttt{HiddenMessage} ghi nhận việc ẩn từng tin nhắn;
\texttt{ConversationMember.historyClearedAt} xác định mốc xóa toàn bộ lịch sử
cũ của thành viên. Khi truy xuất, các tin đã ẩn và tin trước mốc này
không được hiển thị cho người đó; lịch sử của người tham gia khác không
tự động bị xóa. Đây là cách hiểu cần nhóm thống nhất khi chốt thiết kế.

Thuộc tính \texttt{Conversation.isEncrypted} chỉ ghi nhận trạng thái mã hóa
của nhóm. ERD hiện tại chưa mô tả chi tiết dữ liệu phục vụ quản lý khóa,
thay đổi khóa hoặc định dạng nội dung mã hóa; cần hoàn thiện phần này trước
khi khẳng định thiết kế dữ liệu đã đáp ứng đầy đủ chức năng mã hóa nhóm.
Tương tự, chính sách giữ lại hoặc xử lý dữ liệu liên quan khi xóa tài khoản
cần được xác định trước khi chốt các hành vi của khóa ngoại trong script.

\paragraph{Dữ liệu phục vụ thống kê}
\mbox{}\par

Mô hình không tạo bảng thống kê riêng khi số liệu có thể được tính từ
các dữ liệu gốc. Các chức năng của Admin được đối chiếu như sau:

\begin{itemize}
  \item \hyperlink{F14}{F14}: lấy các lượt đăng nhập từ
    \texttt{LoginHistory}, kết hợp hồ sơ tài khoản và sắp xếp theo thời gian.
  \item \hyperlink{F17}{F17}--\hyperlink{F18}{F18}: lọc
    \texttt{UserAccount.createdAt} theo năm đã chọn; tổng hợp số tài khoản
    tạo mới theo tháng của năm được chọn.
  \item \hyperlink{F19}{F19}: đếm quan hệ trong \texttt{Friendship};
    \textbf{Total network} là hợp của tập bạn trực tiếp và tập bạn của bạn,
    loại chính tài khoản đang xét và chỉ đếm mỗi tài khoản một lần.
    \textbf{Friends of friends only} loại thêm các bạn trực tiếp; người bạn
    chung được xác định từ hai quan hệ bạn bè nối tiếp nhau.
  \item \hyperlink{F20}{F20}: đếm lượt mở ứng dụng từ
    \texttt{AppOpenEvent}; đối chiếu tin nhắn, hội thoại và thành viên để
    xác định số người và số nhóm có trao đổi trong năm đã chọn.
  \item \hyperlink{F21}{F21}: đếm các tài khoản khác nhau có mở ứng dụng
    trong từng tháng, không dùng tổng số lượt mở thay cho số người hoạt động.
\end{itemize}

Lịch sử đăng nhập và lượt mở ứng dụng được tách riêng vì việc mở lại ứng dụng
không nhất thiết là một lần đăng nhập mới. Trạng thái online/offline sẽ cần
được backend xác định từ kết nối thực tế khi triển khai; một bản ghi đăng nhập
cũ không đủ chứng minh người dùng đang online.

Bộ chọn \textbf{Period} áp dụng cùng một năm dương lịch cho các bảng và biểu đồ.
Truy vấn theo năm bao gồm các mốc từ đầu năm đến trước đầu năm kế tiếp theo
múi giờ thống nhất của hệ thống. Biểu đồ luôn có 12 tháng; tổng
\textbf{Active users} của cả năm đếm mỗi tài khoản một lần, kể cả khi tài khoản
có hoạt động trong nhiều tháng.

\paragraph{Dữ liệu cho phần mở rộng quản trị}
\mbox{}\par
ERD hiện tại hỗ trợ trạng thái \texttt{ACTIVE}/\texttt{LOCKED} và các báo cáo
spam. Nhật ký khóa/mở khóa trong nguyên mẫu là phần mở rộng cần lưu thời điểm,
tài khoản bị tác động và quản trị viên thực hiện. Nhật ký hoạt động nhóm có thể
lấy mốc gửi tin nhắn và tham gia/rời nhóm từ các bảng hiện có. Để phân biệt
tự rời với bị xóa khỏi nhóm hoặc ghi nhận việc đổi tên và đổi quyền, cần bổ
sung bản ghi sự kiện trước khi triển khai các mục tương ứng.
